% Generated draft: requires mathematical and source review.
\documentclass[11pt]{article}
\usepackage[margin=25mm]{geometry}
\usepackage{fontspec}
\setmainfont{Latin Modern Roman}
\usepackage{amsmath,amssymb,mathtools,unicode-math}
\setmathfont{Latin Modern Math}
\usepackage{enumitem,needspace,xurl,hyperref}
\setlength{\parindent}{0pt}
\setlength{\parskip}{5pt}
\setlength{\emergencystretch}{3em}
\newcommand{\N}{\mathbb{N}}
\newcommand{\Z}{\mathbb{Z}}
\newcommand{\Q}{\mathbb{Q}}
\newcommand{\R}{\mathbb{R}}
\newcommand{\C}{\mathbb{C}}
\newcommand{\F}{\mathbb{F}}
\newcommand{\A}{\mathbb{A}}
\newcommand{\PP}{\mathbb{P}}
\newcommand{\E}{\mathbb{E}}
\newcommand{\eps}{\varepsilon}
\newcommand{\dd}{\,\mathrm{d}}
\begin{document}
\section*{UnsolvedMath ID: 30002783}
\textbf{Source problem code:} catalogue entry
\subsection{Definitions and mathematical statement}
The supplied record groups several intersection questions but omits the parameter restrictions and some hypotheses needed for a fully quantified statement. The following notation makes the requested conclusions explicit; it does not supply those missing hypotheses.

Let \(\Delta^m\) be an \(m\)-dimensional simplex, and let \(\mathcal F(K)\) denote the collection of nonempty faces of a simplex or convex polytope \(K\). For an integer \(q\ge 2\) and a map \(f:K\to Y\), the \(q\)-fold disjoint-face intersection conclusion means
\[
\exists\,\sigma_1,\ldots,\sigma_q\in\mathcal F(K)\quad
\bigl(\sigma_i\cap\sigma_j=\varnothing\text{ for }i\ne j\bigr)
\quad\text{and}\quad
\bigcap_{i=1}^{q} f(\sigma_i)\ne\varnothing.
\]
For a simplex, disjointness of faces is equivalent to disjointness of their vertex sets. The symbol \(q\) is introduced here only to express the number of faces; the record does not specify that number or its relation to its symbols \(m,r\).

A nonempty space \(X\) is contractible if its identity map is homotopic to a constant map. For a family \(\mathcal G\) of subsets of \(Y\), the contractible-intersection condition is
\[
\text{for every finite nonempty }\mathcal H\subseteq\mathcal G,
\qquad
\bigcap_{A\in\mathcal H} A\ne\varnothing
\ \Longrightarrow\ 
\bigcap_{A\in\mathcal H} A\text{ is contractible}.
\]
A good cover commonly also requires its members to be open and to cover the space under consideration. The record does not state which convention applies to face images, which family of faces is used, or which space is covered. Requiring individual face images to be contractible does not, by itself, impose the displayed condition on their intersections.

The supplied questions can therefore be recorded as the following unresolved statement templates.
\begin{enumerate}
\item In the intended topological Tverberg setting, does every continuous map \(f:\Delta^m\to Y\) satisfying the intended good-cover condition on its face images have the \(q\)-fold disjoint-face intersection conclusion? Separately, does the conclusion hold under the intended contractibility condition on face images? The target \(Y\), admissible parameters, precise families of images, and whether these conditions are alternatives or combined hypotheses are not specified.
\item For the intended continuous maps \(f:\Delta^m\to Y\), is
\[
\bigcap_{\substack{\sigma\in\mathcal F(\Delta^m)\\ \dim\sigma=m-r}}
 f(\sigma)\ne\varnothing?
\]
Here \(0\le r\le m\) makes the displayed family consist of nonempty faces of a valid dimension. This range is a condition for interpreting this template, not a recovered source hypothesis. The meaning of \(r\), its relation to the target dimension, and any additional assumptions on \(f\) are missing.
\item For an \(m\)-dimensional convex polytope \(P\) and an affine map \(f:P\to\R^d\), does the \(q\)-fold disjoint-face intersection conclusion hold under the intended hypotheses? Affine means
\[
f((1-t)x+ty)=(1-t)f(x)+tf(y)
\quad(x,y\in P,\ 0\le t\le1).
\]
The Euclidean target in this template provides a conventional setting for the word affine; the record does not explicitly specify the target, admissible polytopes, or restrictions relating \(m,d,q\).
\end{enumerate}
No universal assertion over unrestricted choices of these parameters is intended. The missing specifications prevent recovery of a unique precise conjecture from the supplied record.
\subsection{Short English statement}
The record asks whether Tverberg intersections persist under good-cover or contractibility assumptions, whether all images of faces of dimension \(m-r\) share a point, and whether an analogous disjoint-face conclusion holds for affine maps on polytopes. Essential parameter restrictions and hypotheses are absent, so these are explicit statement templates requiring human mathematical review. The supplied partially solved label is dated curation, not an independently verified conclusion.
\subsection{Sources}
\begin{itemize}
\item Local registry: \texttt{data/unsolvedmath/problems.json} and supplementary registry \texttt{catalog/problems.json}; supplied permanent UnsolvedMath ID \(30002783\), source problem code \texttt{catalogue entry}, title \emph{Tverberg Intersections Under Good Cover Conditions}. These registry references identify the supplied record; their contents were not independently inspected.
\item Supplied report citation: \emph{Geometric and Algebraic Combinatorics} (2015), Oberwolfach Reports, pp.\ 285--368, DOI 10.4171/owr/2015/5. \url{https://doi.org/10.4171/owr/2015/5}. The background describes Oberwolfach Reports as documenting mathematical workshops and their open-problem discussions. No individual proposer or exact problem page is supplied; the proposed year is 2015.
\item Supplied literature source: Frick, counterexamples to topological Tverberg. \url{https://arxiv.org/abs/1502.00947}.
\item Supplied literature source: Karasev, topological central point theorem. \url{https://doi.org/10.1016/j.topol.2011.12.002}.
\item Supplied literature source: original OWR report. \url{https://doi.org/10.4171/owr/2015/5}.
\item Supplied literature assessment, checked \texttt{2026-08-21}: the registry labels the record partially solved and states that the affine Tverberg theorem, prime-power topological cases, and the stated central-point consequence are established; that Frick's counterexamples disprove unrestricted continuous TTC beyond prime powers; and that the good-cover variants remain unresolved. This is supplied dated curation derived from the OWR record, not independently verified literature assessment. Its time-sensitive status and cited claims require independent verification, and its shorthand does not recover the omitted hypotheses.
\item The registry labels the statement \texttt{corrected\_verified} and reports that notation and publication prose were repaired without changing its substance. This is a supplied editorial qualification, not independent verification of the mathematical formulation or of the linked sources.
\end{itemize}
\end{document}
